\documentclass[fleqn]{report}

\usepackage[serbian]{babel}
\usepackage{csquotes}
\usepackage[a4paper]{geometry}
\usepackage{amsmath}
\usepackage{amssymb}
\usepackage{siunitx}
\usepackage{cancel}
\usepackage{empheq}
\usepackage{tcolorbox}

\sisetup{
    per-mode = symbol
}

\tcbuselibrary{breakable}
\newtcolorbox{tbox} {
    breakable,
    colback=gray!5,
    colframe=gray!60,
    arc=2mm,
    boxrule=0.5pt,
    top=1mm, bottom=1mm, left=3mm, right=3mm
}

\title{
    Laserska tehnika \\
    \large bilješke sa predavanja
}
\author{Luka Ristanović}
\date{Školska 2026/27. godina}

\begin{document}
\maketitle

\tableofcontents

\chapter{Uvod}

\section{Bezbjednost pri radu sa laserima}

Pri radu sa laserima bitno je voditi računa o sljedećim mogućim opasnostima:
\begin{enumerate}
    \item električni udar
          \subitem U propratnoj elektronici koja ide uz lasere, primarno napajanje, često se pojavljuju kondenzatori, koji mogu ostati napunjeni čak i kada je napajanje isključeno, što predstavlja bezbjednosni rizik, stoga ih je neophodno isprazniti i u tom slučaju.
          Obavezno je pretpostaviti da laser nije bezbjedan čak i kad izgleda pouzdano i atraktivno ("fancy"). Ne smije se dirati ništa dok laser radi.
    \item oštećenja oka
          \subitem Za ljudski vid najopasnije su talasne dužine 400um - 1400um!
          Poluprovodnički laser od 15mW koji radi na 870um trajno oštećuje retinu.
          Voditi računa i o refleksijama!
          Za talasne dužine koje se ne vide obavezne su naočare!
          Za vidljive talasne dužine ljudi teže da skidaju naočare, što je nebezbjedno.
          U krajnjem slučaju može doći i do odvajanja retine.
    \item trovanje
          \subitem Hemikalije za laserske diode su veoma otrovne.
          Neki kristali za nelinearnu optiku su takođe otrovni i (?).
    \item opekotine
          \subitem $CO_2$ laser može opeći kožu.
          Talasna dužina od 276um oštećuje DNK, što je naročito opasno jer je to drugi harmonik talasne dužine od 552um, na kojoj radi značajan broj lasera.
\end{enumerate}

\section{Osnovni pojmovi}

Upoznajmo se sada sa nekim od osnovnih pojmova koji se koriste u laserskoj tehnici.
Svaka tema će naknadno biti detaljnije obrađena, pa ćemo sada preći samo najosnovnije.

\subsection{Atomske orbitale i energetska stanja elektrona}

Elektroni u atomu mogu postojati u dobro definisanim energetskim ljuskama opisanim glavnim kvantnim brojem, pri čemu energija elektrona koji se nalazi u jednoj ljusci sa tim brojem raste.
Unutar tih glavnih ljuski elektroni mogu zauzimati tzv. elektronske orbitale (podljuske), koje imaju oznake redom: s(2), p(6), d(10), f(14), g(18)... koje takođe imaju rastuće energije, tako da svaka orbitala u svakoj ljusci ima različitu odgovarajuću energiju (broj u zagradi je broj elektrona koji mogu da se nalaze u datoj orbitali).
U ravnoteži, elektroni u atomu redom popunjavaju stanja najnižih energija, pa je elektronska struktura atoma jednoznačno određena orbitalama koje su tada popunjene. Na primjer, elektronsku strukturu kriptona možemo zapisati na sljedeći način:

\begin{equation*}
    \text{Kr: } 1s^2 2s^2 2p^6 3s^2 3p^6 3d^{10} 4s^2 4p^6
\end{equation*}

Pojedinim procesima, moguće je izbaciti atom iz pomenute ravnoteže, odnosno predati energiju nekom elektronu, tako da on pređe na neko slobodno stanje unutar atoma sa višom energijom, a prethodno sa nižom ostavi upražnjenim.
Pomenuti prelazi se u laserima najčešće dešavaju između najudaljenijih ljuski u atomu (sa najvišom energijom).
Takođe se mogu dešavati uvijek kada postoje energetska stanja koja takve prelaze dozvoljavaju, bez obzira na to da li je u pitanju čvrsto tijelo, tečnost, gas, plazma.
Ta stanja takođe mogu biti i dva vibraciona moda molekula (npr. $CO_2$), ili valentna i provodna zona u poluprovodniku.

\subsection{Inverzna naseljenost nivoa}

Kao što smo rekli, u ravnoteži elektroni teže da popunjavaju energetska stanja nižih energija, što je opisano Bolcmanovom (Gibsovom) raspodjelom:
\begin{equation*}
    N_n \propto e^{-\frac{E_n}{k_B T}}
\end{equation*}
gdje je $N_n$ koncentracija elektrona na $n$-tom stanju, a $E_n$ energija tog stanja.
Odatle slijedi da je odnos koncentracija elektrona na stanjima 1 i 2 jednak:
\begin{equation*}
    \frac{N_2}{N_1} = e^{-\frac{E_2 - E_1}{k_B T}}
\end{equation*}
Da bi došlo do takozvanog laserovanja, kao što ćemo pokazati, mora biti ostvarena \emph{inverzna naseljenost nivoa}, odnosno moraju postojati stanja u kojima ona sa višom energijom imaju višu koncentraciju elektrona:
\begin{equation*}
    N_2 > N_1
\end{equation*}
Na prvi pogled, izgleda kao da temperatura mora biti negativna da bi ovo bilo ostvarivo.
Međutim, Bolcmanova raspodjela važi za termodinamičku ravnotežu.
Stanje sa višom energijom moguće je zato \enquote{pumpati} elektronima nekim neravnotežnim procesom.
Za to se mogu koristiti svjetlost, električno pražnjenje, hemijski procesi itd.

\subsection{Spontana i stimulisana emisija}

Nakon što se nekim procesom pumpanja elektroni dovedu na više energetsko stanje, oni mogu preći nazad jednim od dva procesa koji su nama od interesa:

\begin{enumerate}
    \item \underline{spontana emisija} --- proces u kom elektron spontano prelazi sa višeg na niže energetsko stanje pri čemu se razlika u energijama oslobađa u vidu fotona energije takve da je
          \begin{equation*}
              E_{ph} = h \nu = E_2 - E_1
          \end{equation*}
          Emitovani foton osim te učestanosti ima potpuno nasumične karakteristike (pravac, fazu i polarizaciju), tako da više uzastopnih spontanih emisija proizvodi potpuno nekorelisane fotone, koji čine takozvanu \emph{nekoherentnu} svjetlost.
          Proces spontane emisije najčešće se opisuje \emph{vremenom spontane emisije}, koje predstavlja srednje vrijeme koje elektron, nakon pobuđivanja, provede na višem energetskom stanju, prije nego što se spontano spusti na niže.
    \item \underline{stimulisana emisija} --- proces u kom foton  nailazi na elektron u višem energetskom stanju, interaguje s njim, i \enquote{stimuliše} ga da pređe na niže energetsko stanje.
          U ovom slučaju se pri tom prelazu takođe emituje još jedan foton, ali ovaj put on ima identične karakteristike kao i nailazeći (fazu, polarizaciju i pravac).
\end{enumerate}
Napomenimo odmah, što će kasnije biti od značaja, da u svim ovim procesima, kao i obično, mora važiti zakon održanja impulsa.

\subsection{Snaga i energija}

U opštem slučaju, laseri mogu da rade u dva režima: kontinualni režim - CW (engl. \emph{continuous working}) i impulsni režim (za koji će biti vezano Q-prekidanje i zaključavanje modova).
Svi laseri koji mogu da rade u kontinualnom režimu, mogu da rade i u impulsnom. Međutim, obrnuto ne važi, jer pojedini laseri ne mogu da obezbijede dovoljno pumpanja da nadomjeste kontinualnu potrošnju za generisanje zračenja.

Takođe, dodajmo da generalno laseri imaju male efikasnosti konverzije električne energije u optičku.

\subsection{Monohromatičnost, koherencija i širina linije}

Osobine laserskog zračenja, u poređenju sa ostalim izvorima, su usmjerenost, velika iradijansa $\left[\frac{W}{m^2}\right]$ u maloj zapremini prostora i koherentnost (šta god da je to tačno).


\chapter{Gausov snop}
Neki uvod ovdje bro

\section{Paraksijalna Helmholcova jednačina}

Pođimo od Maksvelovih jednačina u dielektriku ili vakuumu:
\begin{align}
     & \nabla \cdot \vec D = \nabla \cdot (\varepsilon \vec E) = \rho = 0                                        & (\rho = 0) \nonumber                \\
     & \nabla \times \vec H = \varepsilon \frac{\partial \vec E}{\partial t}                                     & (\vec J = 0) \label{gs - faradejev} \\
     & \nabla \times \vec E = - \mu \frac{\partial \vec H}{\partial t} \bigg/ \nabla \times \label{gs - amperov}
\end{align}
\begin{align*}
     & \nabla \times (\nabla \times \vec E) = \nabla \times \left(-\mu \frac{\partial \vec H}{\partial t}\right) \\
     & \nabla (\nabla \vec E) - \nabla^2 \vec E = - \mu \frac{\partial}{\partial t} (\nabla\times\vec H)
\end{align*}
\begin{tbox}
    \begin{align*}
         & 0 = \nabla\vec D = \nabla(\varepsilon\vec E) = \nabla\varepsilon \vec E + \varepsilon\nabla\vec E \\
         & \implies \nabla\vec E = -\frac{1}{\varepsilon} \nabla\varepsilon \vec E
    \end{align*}
\end{tbox}
\begin{align*}
     & \nabla\left(-\frac{1}{\varepsilon}\nabla\varepsilon\vec E \right) - \Delta \vec E = -\mu\frac{\partial}{\partial t}\left(\varepsilon\frac{\partial\vec E}{\partial t}\right) \\
     & \Delta\vec E - \mu\varepsilon\frac{\partial^2\vec E}{\partial t^2} = \nabla\left(-\frac{1}{\varepsilon}\nabla\varepsilon\vec E\right)
\end{align*}
Ako se radi o lokalno homogenim sredinama, tada se $\varepsilon(\vec r)$ malo mijenja u okviru jedne talasne dužine, pa možemo aproksimirati da je $\frac{\nabla\varepsilon}{\varepsilon} \approx 0$:
\begin{equation*}
    \Delta\vec E - \mu\varepsilon\frac{\partial^2\vec E}{\partial t^2} = 0
\end{equation*}
Prepostavimo da je električno polje prostoperiodičnog oblika $\vec E(\vec r, t) = \vec E(\vec r)e^{-j\omega t}$:
\begin{align*}
     & \frac{\partial^2 \vec E}{\partial t^2} = - \vec E(\vec r) \omega^2 e^{-j\omega t} = \vec E(\vec r, t) (-\omega^2) \\
     & \implies \Delta\vec E + \omega^2\mu\varepsilon(\vec r)\vec E(\vec r, t) = 0
\end{align*}
Ako uvedemo parametar $k^2(\vec r) = \omega^2\mu\varepsilon(\vec r)$ dobijamo:
\begin{equation*}
    \Delta\vec E + k^2(\vec r)\vec E = 0
\end{equation*}
Ovo predstavlja \emph{Helmholcovu jednačinu}, odnosno talasnu jednačinu za vektor električnog polja u dielektriku u kom nema gubitaka ili pojačanja.
\begin{tbox}
    Ukoliko želimo da uključimo i te efekte, možemo uvesti efektivnu dielektričnu permitivnost kao kompleksnu veličinu:
    \begin{equation*}
        \underline{\varepsilon}(\vec r) = \varepsilon(\vec r)\left(1 - \frac{j\sigma(\vec r)}{\omega\varepsilon}\right)
    \end{equation*}
\end{tbox}
Od interesa nam je da analiziramo prostiranje svjetlosnog talasa koje se formira u laseru i izlazi iz njega.
U praksi, poznato nam je da je takva svjetlost fokusirana u neku vrstu snopa, tako da propagira dominantno u jednom pravcu, dok se u transverzalnom gubi brzo.
Pokušajmo iz tog razloga da riješimo ovu jednačinu u cilindričnom koordinatnom sistemu.
U tom slučaju laplasijan možemo zapisati kao $\Delta = \Delta_t + \frac{\partial^2}{\partial z^2}$, gdje je $\Delta_t = \frac{1}{\rho}\frac{\partial}{\partial \rho} \left( \rho\frac{\partial}{\partial\rho} \right) + \frac{1}{\rho^2}\frac{\partial^2}{\partial\phi^2}$ transverzalni laplasijan.

Dalje, pošto pretpostavljamo da se talas prostire uglavnom po $z$-osi (odnosno da se energija sprovodi u tom pravcu), možemo uvesti \emph{paraksijalnu aproksimaciju} $k \approx k_z$ i pretpostaviti rješenje u obliku $E(\rho, \phi, z) = \Psi(\rho, \phi, z) e^{-jkz}$, gdje je $\Psi$ kompleksna anvelopa talasa:
\begin{align*}
     & \Delta_t E + \frac{\partial^2}{\partial z^2}E + k^2 E = 0                                                                                                                                                        \\
     & \Delta_t\Psi \cancel{e^{-jkz}} + \frac{\partial^2\Psi}{\partial z^2} \cancel{e^{-jkz}} + 2 \frac{\partial\Psi}{\partial z} (-jk) \cancel{e^{-jkz}} + \Psi(-k^2)\cancel{e^{-jkz}} + k^2\Psi \cancel{e^{-jkz}} = 0 \\
     & \Delta_t\Psi + \frac{\partial^2\Psi}{\partial z^2} - 2jk \frac{\partial\Psi}{\partial z} - \cancel{k^2\Psi} + \cancel{k^2\Psi} = 0                                                                               \\
     & \Delta_t\Psi + \frac{\partial^2\Psi}{\partial z^2} - 2jk \frac{\partial\Psi}{\partial z} = 0
\end{align*}
Ako pretpostavimo da je varijacija anvelope u $z$ pravcu dovoljno spora $k\frac{\partial\Psi}{\partial z} \gg \frac{\partial^2\Psi}{\partial z^2}$, možemo zanemariti drugi izvod u gornjem izrazu i ostaje:
\begin{align}
     & \Delta_t\Psi - 2jk\frac{\partial\Psi}{\partial z} = 0 \label{gs - phj}                                                                                                                                       \\
     & \frac{1}{\rho}\frac{\partial}{\partial \rho} \left( \rho\frac{\partial\Psi}{\partial\rho} \right) + \frac{1}{\rho^2}\frac{\partial^2\Psi}{\partial\phi^2} - 2jk\frac{\partial\Psi}{\partial z} = 0 \nonumber
\end{align}
Takođe pretpostavljamo da nema varijacije po azimutnom uglu $\Psi \neq \Psi(\phi) \implies \frac{\partial^2\Psi}{\partial\phi^2} = 0$:
\begin{equation*}
    \frac{1}{\rho}\frac{\partial}{\partial\rho} \left( \rho\frac{\partial\Psi}{\partial\rho}\right) - 2jk\frac{\partial\Psi}{\partial z} = 0
\end{equation*}
Ovo nazivamo \emph{paraksijalnom Helmholcovom jednačinom} (PHJ) za kružnu aperturu.

Rješenje svake diferencijalne jednačine je svaka funkcija koja zadovoljava tu jednačinu.
Da bi rješenje bilo fizičko i predstavljalo stvarnost, zadovoljavanje jednačine je potreban uslov, ali ne i dovoljan.

Jedno rješenje je, na primjer, parabolični talas $\Psi(\rho, z) = \frac{\Psi_0}{z}e^{-jk\frac{\rho^2}{2z}}$, što je rješenje najnižeg reda:
\begin{align*}
    \frac{\partial\Psi}{\partial z} & = \frac{\Psi_0}{z^2}\left(ze^{-jk\frac{\rho^2}{2z}}(-jk\rho^2)(-\frac{1}{2z^2}) - e^{-jk\frac{\rho^2}{2z}}\right) = \\
                                    & = \Psi_0e^{-jk\frac{\rho^2}{2z}}\left(\frac{jk\rho^2}{2z^3} - \frac{1}{z^2}\right)
\end{align*}
\begin{align*}
     & \rho\frac{\partial\Psi}{\partial\rho} = \rho\frac{\Psi_0}{z}e^{-jk\frac{\rho^2}{2z}}(-jk)\frac{\rho}{z} = -jk\frac{\Psi_0}{z^2}\rho^2e^{-jk\frac{\rho^2}{2z}}                                                                                                                                 \\
     & \frac{\partial}{\partial\rho}\left(\rho\frac{\partial\Psi}{\partial\rho}\right) = -jk\frac{\Psi_0}{z^2}\left(2\rho e^{-jk\frac{\rho^2}{2z}} + \rho^2e^{-jk\frac{\rho^2}{2z}}(-\frac{jk\rho}{z})\right) = -jk\frac{\Psi_0}{z^2}\left(2\rho - \frac{jk\rho^3}{z}\right)e^{-jk\frac{\rho^2}{2z}} \\
     & \frac{1}{\rho}\frac{\partial}{\partial\rho}\left(\rho\frac{\partial\Psi}{\partial\rho}\right) = -2jk\frac{\Psi_0}{z^2} \left(1 - \frac{jk\rho^2}{2z}\right)e^{-jk\frac{\rho^2}{2z}} = 2jk\Psi_0\left(\frac{jk\rho^2}{2z^3} - \frac{1}{z^2} \right)e^{-jk\frac{\rho^2}{2z}}
\end{align*}
Vidimo da za ovakvo $\Psi$ zaista važi $\frac{1}{\rho}\frac{\partial}{\partial\rho}\left(\rho\frac{\partial\Psi}{\partial\rho}\right) - 2jk\frac{\partial\Psi}{\partial z} = 0$.
\begin{tbox}
    Parabolični talas je paraksijalna aproksimacija sfernog talasa, dakle kada uzmemo da su $x$ i $y$ mnogo manji od $z$.
    Krenemo od jednačine sfernog talasa:
    \begin{equation*}
        \Psi(\vec r) = \frac{A_1}{r} e^{-jkr}
    \end{equation*}
    i uvedimo Frenelovu aproksimaciju paraboličnog talasa $\sqrt{x^2 + y^2} \ll z$:
    \begin{equation*}
        r = \sqrt{x^2 + y^2 + z^2} = z\sqrt{1 + \frac{x^2 + y^2}{z^2}} = z\left(1 + \frac{x^2 + y^2}{z^2}\right)^{\frac{1}{2}} \approx z \left(1 + \frac{1}{2}\frac{x^2 + y^2}{z^2}\right)
    \end{equation*}
    Znamo da je $x^2 + y^2 = \rho^2$ pa je $r^2 \approx z(1 + \frac{\rho^2}{2z^2}) = z + \frac{\rho^2}{2z}$ i:
    \begin{equation*}
        \Psi(\vec r) = \frac{A_1}{z} e^{-jkz} e^{-jk\frac{\rho^2}{2z}}
    \end{equation*}
    \emph{Napomena: $r$ u imeniocu smo jednostavno aproksimirali sa $z$, dok u eksponentu to ne smijemo da uradimo zato što se $\frac{\rho^2}{2z}$ množi sa $k$ koje može biti veliko, pa taj proizvod može prouzrokovati značajne promjene u fazi.}
\end{tbox}
Problem sa ovim konkretnim rješenjem, iako zadovoljava samu jednačinu, je taj što je nefizičko: za $z \rightarrow 0$ dobijamo da amplituda polja eksplodira u beskonačnost, i ovakav talas nema jasno formiran, konfiniran, uzak snop.

Kako riješiti ovaj problem? Prvo, konstatujmo da se u PHJ $z$ koordinata javlja samo u okviru izvoda, odnosno ta jednačina predstavlja linearnu diferencijalnu jednačinu sa konstantnim koeficijentima u odnosu na $z$.
To znači da, ako uvedemo neku novu promjenljivu $q = z - \xi$, gdje je $\xi$ proizvoljna konstanta, jednačinu možemo transformisati na sledeći način:
\begin{align*}
     & \Delta_t\Psi - 2jk\frac{\partial\Psi}{\partial z} = 0                            \\
     & \Delta_t\Psi - 2jk\frac{\partial\Psi}{\partial q}\cancelto{1}{\frac{dq}{dz}} = 0
\end{align*}
Dobili smo diferencijalnu jednačinu po $\Psi(\rho, q)$ koja ima identičnu formu kao jednačina \eqref{gs - phj} po $\Psi(\rho, z)$.
Odavde slijedi da rješenja za ovu jednačinu $\Psi(\rho, q)$ imaju identičan oblik kao i rješenja za \eqref{gs - phj}, dakle možemo napisati da je:
\begin{equation}
    \Psi(\rho, q) = \frac{\Psi_0}{q}e^{-jk\frac{\rho^2}{2q}} \label{gs - psi(rho, q)}
\end{equation}
pri čemu je $q = q(z) = z - \xi$, za proizvoljnu konstantu $\xi$.
Konačno, da bismo formirali rješenje koje je od fizičkog značaja i našeg interesa, pogledajmo šta se dešava ako izaberemo $\xi$ tako da bude čisto imaginarno: $\xi = -jz_0$, tako da je $z_0$ realna, pozitivna konstanta.
Dakle, možemo da pišemo $q(z) = z - \xi = z + jz_0$.
Konkretnije, iz razloga koji će postati jasni naknadno (da bismo iz samog izraza mogli što lakše da utvrdimo pojedine željene parametre), napišimo $q(z)$ preko novih promjenljivih $R(z)$ i $W(z)$ na sljedeći način:
\begin{equation*}
    \frac{1}{q(z)} = \frac{1}{R(z)} - j\frac{\lambda}{\pi W^2(z)}
\end{equation*}
gdje su:
\begin{empheq}[box=\fbox]{align*}
    & R(z) = z\left(1 + \left(\frac{z_0}{z}\right)^2\right) \\
    & W(z) = W_0 \sqrt{1 + \left(\frac{z}{z_0}\right)^2}    \\
    & W_0 = \sqrt{\frac{\lambda z_0}{\pi}}                  \\
\end{empheq}
\begin{tbox}
    Provjerimo da li je ovaj izraz zaista tačan:
    \begin{align*}
        \frac{1}{R(z)} - j\frac{\lambda}{\pi W^2(z)} & = \frac{1}{z\left(1 + \left(\frac{z_0}{z}\right)^2\right)} - j\frac{\lambda}{\pi W_0^2\left(1 + \left(\frac{z}{z_0}\right)^2\right)} =                    \\
                                                     & = \frac{z}{z^2 + z_0^2} - j\frac{\cancel{\lambda}}{\cancel{\pi} \frac{\cancel{\lambda} z_0}{\cancel{\pi}}\left(1 + \left(\frac{z}{z_0}\right)^2\right)} = \\
                                                     & = \frac{z}{z^2 + z_0^2} - j\frac{z_0}{z_0^2 + z^2} =                                                                                                      \\
                                                     & = \frac{z - jz_0}{z^2 + z_0^2} =                                                                                                                          \\
                                                     & = \frac{\cancel{z - jz_0}}{\cancel{(z - jz_0)}(z + jz_0)} =                                                                                               \\
                                                     & = \frac{1}{z + jz_0} =                                                                                                                                    \\
                                                     & = \frac{1}{q(z)}
    \end{align*}
\end{tbox}
\begin{tbox}
    Veličina $z_0 = \frac{\pi W_0^2}{\lambda}$ se naziva Rejlijevo rastojanje.
\end{tbox}
Ako ovakav izraz za $q(z)$ ubacimo u izraz \eqref{gs - psi(rho, q)}, dobijamo:
\begin{align*}
    \Psi(\rho, q(z)) = \Psi(\rho, z) & = \Psi_0 \frac{1}{z + jz_0} e^{-\frac{jk\rho^2}{2}\left(\frac{1}{R(z)} - j\frac{\lambda}{\pi W(z)^2}\right)} =                                                                                                 \\
                                     & = \Psi_0 \frac{1}{jz_0}\frac{1}{1 - j\frac{z}{z_0}} e^{-\frac{jk\rho^2}{2}\left(\frac{1}{R(z)} - j\frac{\lambda}{\pi W(z)^2}\right)} =                                                                         \\
                                     & = \frac{\Psi_0}{jz_0} \frac{1}{\sqrt{1 + \left(\frac{z}{z_0}\right)^2}}e^{-j\arctg\frac{z}{z_0}} e^{-\frac{jk\rho^2}{2R(z)}} e^{-\frac{\cancel{k\lambda}\rho^2}{\cancel{2\pi} W(z)^2}} =                       \\
                                     & = \underbrace{\frac{\Psi_0}{jz_0}}_{\Psi_0'} \underbrace{\frac{1}{\sqrt{1 + \left(\frac{z}{z_0}\right)^2}}}_{\frac{W_0}{W(z)}} e^{-\frac{\rho^2}{W(z)^2}} e^{-jk\frac{\rho^2}{2R(z)} - j\arctg\frac{z}{z_0}} = \\
                                     & = \Psi_0' \frac{W_0}{W(z)} e^{-\frac{\rho^2}{W(z)^2}} e^{-jk\frac{\rho^2}{2R(z)} - j\arctg\frac{z}{z_0}}
\end{align*}
Onda je kompleksna amplituda električnog polja:
\begin{empheq}[box=\fbox]{align}
    E(\rho, z) & = \Psi(\rho, z) e^{-jkz} = \nonumber                                                                                     \\
    & = \Psi_0' \frac{W_0}{W(z)} e^{-\frac{\rho^2}{W(z)^2}} e^{-jk\frac{\rho^2}{2R(z)} - j\arctg\frac{z}{z_0} - jkz} \label{gs - E}
\end{empheq}
Znamo da je iradijansa $I$ proporcionalna sa kvadratom modula električnog polja $|E|^2 = EE^*$:
\begin{empheq}[box=\fbox]{equation*}
    I = I_0 \left(\frac{W_0}{W(z)}\right)^2 e^{-\frac{2\rho^2}{W(z)^2}}
\end{empheq}
Ovi izrazi za električno polje i iradijansu konačno opisuju Gausov snop.

\section{Osobine Gausovog snopa}

Analizirajmo detaljno sada izraz za iradijansu gausovog snopa $I(\rho, z)$.
Prvo, vidimo da u jednoj ravni normalnoj na pravac snopa (dakle za $z = const$), zavisnost iradijanse od radijalne koordinate ima oblik Gausove zvonaste krive, sa pikom na samoj osi intenziteta $I_0\left(\frac{W_0}{W(z)}\right)^2$.

Dalje, ako posmatramo baš taj intenzitet na osi snopa u zavisnosti od $z$, nakon što ubacimo da je $W(z) = W_0 \sqrt{1 + \left(\frac{z}{z_0}\right)^2}$ dobijamo:
\begin{equation*}
    I(0, z) = I_0 \frac{1}{1 + \left(\frac{z}{z_0}\right)^2}
\end{equation*}
Vidimo da intenzitet pika iradijanse (na osi snopa) opada sa rastojanjem od izvora, odnosno funkcija $I(\rho, z=const)$ se \enquote{spljoštava}.

(Grafici I(ro) za razlicite z )

Odredimo sada pojedine parametre od interesa za ovakav snop.

\subsection{Snaga}
Poznato je da je snaga $P = \int_S IdS$, gdje se integrali po nekoj površi kroz koju se talas prostire, normalno na njegov pravac.
Talas se prostire u $z$ pravcu, pa možemo integraliti po ravni $z=const$:
\begin{align*}
    P & = \int_{\rho=0}^\infty \int_{\phi = 0}^{2\pi} I \underbrace{\rho d\rho d\phi}_{dS} =                                                                           \\
      & = \underbrace{\int_0^{2\pi}d\phi}_{2\pi} \int_0^\infty I_0\left(\frac{W_0}{W(z)}\right)^2 \rho e^{-\frac{2\rho^2}{W(z)^2}} d\rho =                             \\
      & = 2\pi I_0 \left(\frac{W_0}{W(z)}\right)^2 \int_0^\infty \rho e^{-\frac{2\rho^2}{W(z)^2}} =                                                                    \\
      & = 2\pi I_0 \left(\frac{W_0}{\cancel{W(z)}}\right)^2 \frac{\cancel{W(z)^2}}{4} \int_0^\infty e^{-\frac{2\rho^2}{W(z)^2}} d\left(\frac{2\rho^2}{W(z)^2}\right) = \\
      & = \frac{\pi}{2} I_0 W_0^2 \underbrace{\left.e^{-\frac{2\rho^2}{W(z)^2}}\right|_\infty^0}_1 =                                                                   \\
      & = \frac{\pi}{2} I_0 W_0^2
\end{align*}

Laserski snop se često karakteriše i svojom snagom $P$, pa može biti korisno zapisati i iradijansu preko nje na sljedeći način:
\begin{equation*}
    I(\rho, z) = \frac{2P}{\pi W(z)^2} e^{-\frac{2\rho^2}{W(z)^2}}
\end{equation*}

Još jedna stvar koja može biti od interesa, nasuprot ukupne snage snopa, je procenat njegove ukupne snage koji se prostire unutar radijusa $\rho_0$ od ose:
\begin{equation*}
    \frac{1}{P} \int_0^{\rho_0} I(\rho, z) 2\pi\rho d\rho = 1 - e^{-\frac{2\rho_0^2}{W(z)^2}}
\end{equation*}
Primjera radi, dobija se da se približno 86\% snage prostire unutar $\rho_0 = W(z)$, a 99\% snage unutar $\rho_0 = 1.5W(z)$.

\subsection{Poluprečnik snopa}
Gausov snop nema egzaktan poluprečnik, u smislu radijalnog rastojanja na kom polje opadne na nulu, već od zavisnosti od njega ono opada eksponencijalno.
Zbog toga se poluprečnik snopa definiše kao radijalno rastojanje na kom vrijednost amplitude polja opadne $e$ puta od vrijednosti na osi, odnosno iradijansa opadne $e^2$ puta:
\begin{align*}
     & \frac{I(\rho, z)}{I(0, z)} = e^-2                                                                                                          \\
     & \frac{\cancel{I_0 \left(\frac{W_0}{W(z)}\right)^2} e^{-\frac{2\rho^2}{W(z)^2}}}{\cancel{I_0 \left(\frac{W_0}{W(z)}\right)^2} e^0} = e^{-2} \\
     & \frac{2\rho^2}{W(z)^2} = 2                                                                                                                 \\
     & \rho = W(z)
\end{align*}
Dakle poluprečnik iznosi $W(z) = W_0\sqrt{1 + \left(\frac{z}{z_0}\right)^2}$.
Minimalna vrijednost poluprečnik dostiže za $z = 0$ i iznosi:
\begin{equation*}
    W_{min} = W(z = 0) = W_0
\end{equation*}
i naziva se \emph{poluprečnik struka snopa} (pa je i $z = 0$ položaj struka snopa).
Prečnik struka snopa $2W_0$ naziva se još i veličina spota (engl. \emph{spot size}).

(Grafik W(z))

\subsection{Divergencija snopa}

Sa grafika $W(z)$ možemo uočiti da za $z \gg z_0$ funkcija poprima linearan oblik:
\begin{equation*}
    W(z) = W_0 \sqrt{1 + \left(\frac{z}{z_0}\right)^2} \approx W_0 \sqrt{\left(\frac{z}{z_0}\right)^2} = \frac{W_0}{z_0}z
\end{equation*}
odnosno, nakon što se dovoljno udalji od struka, snop se u odnosu na osu širi, tj. divergira, pod uglom $\theta_0$ tako da je:
\begin{align*}
    \theta_0 \approx \tg\theta_0 = \frac{W_0}{z_0} = \frac{W_0}{\frac{\pi W_0}{\lambda}} = \frac{\lambda}{\pi W_0} \text{($\theta$ je u praksi dovoljno malo da možemo pisati $\theta \approx \tg\theta$)}
\end{align*}
Ovaj ugao definiše \emph{divergenciju snopa}.
Vidimo da je ona obrnuto srazmjerna sa širinom struka $W_0$, zbog čega je teško formirati uzak snop na velikom rastojanju (ako izaberemo malu širinu struka, snop će brzo divergirati i raširiti se, a ako izaberemo veliku širinu struka, snop je onda već inicijalno širok i odatle se samo još širi).

Primijetimo još da važi:
\begin{equation*}
    \theta_0 W_0 = \frac{\lambda}{\pi} = const
\end{equation*}

\subsection{Dubina fokusa}

\emph{Dubina fokusa} ili \emph{konfokalni parametar} $D$ predstavlja mjeru divergencije snopa.
Definiše se kao raspon $z$ koordinate za koju se površina snopa udvostruči.
\begin{align*}
     & \cancel{\pi} W(z)^2 = 2 \cancel{\pi} W_0^2                                     \\
     & \cancel{W_0^2} \left(1 + \left(\frac{z}{z_0}\right)^2\right) = 2\cancel{W_0^2} \\
     & \left(\frac{z}{z_0}\right)^2 = 1                                               \\
     & z = z_0                                                                        \\
     & \implies D = 2z = 2z_0 = \frac{2\pi W_0^2}{\lambda}
\end{align*}
Od interesa je da imamo što usmjereniji snop koji će preći što veće rastojanje prije nego što mu se površina udvostruči, dakle želimo što veće $D$.
Da bi to bilo ispunjeno, očigledno je potrebno da imamo što širi struk snopa $W_0$ i što manju talasnu dužinu $\lambda$.

Primijetimo odmah da je dubina fokusa proporcionalna površini spota $D \propto \pi W_0^2$.
Ovo uvodi jedan značajan \enquote{trade-off} u laserskoj tehnici koji je sledeći:
ako imamo neki laserski snop koji je širi nego što je nama potrebno, možemo određenom aparaturom (na primjer uz pomoć sočiva) da ga modifikujemo i formiramo mu još jedan struk na mjestu gdje je nama potrebno da bude fokusiran.
Problem nastaje ako je nama potrebno da snop bude fokusiran na veoma malu površinu, zato što tada onda i dubina fokusa mora biti mala.
Samim tim je i mali prostor za grešku tokom pozicioniranja struka ili objekta na koji želimo da fokusiramo lasersko zračenje, što se može pokazati kao inženjerski izazovnim.
Na primjer, za HeNe laser, koji radi na talasnoj dužini $\lambda = 633\si{\milli\meter}$, prečniku struka $2W_0 = 2\si{\centi\meter}$ odgovara dubina fokusa $D \approx 1\si{\kilo\meter}$, a $2W_0 = 20\si{\micro\meter}$ odgovara $D \approx 1\si{\milli\meter}$.

\subsection{Faza talasa}
Iz izraza za električno polje Gausovog snopa vidimo da je faza talasa:
\begin{equation*}
    \phi(\rho, z) = kz - \xi(z) + \frac{k\rho^2}{2R(z)}
\end{equation*}
\begin{tbox}
    Znak minus ispred $\xi(z)$ potiče od promjene faze za $\frac{\pi}{2}$ usljed \enquote{upijanja} $j$ u $\Psi_0'$.
\end{tbox}
Član $kz$ pokazuje sličnost Gausovog snopa sa ravanskim talasom, član $\frac{k\rho^2}{2R(z)}$ frontovima daje zakrivljenost, a $\xi(z) = \arctg\frac{z}{z_0} = \arctg\frac{z\lambda}{\pi W_0^2}$ je odstupanje.

(Grafik za ksi)

Specifično, ako posmatramo fazu talasa na osi snopa, vidimo da je:
\begin{equation*}
    \phi(0, z) = kz - \xi(z)
\end{equation*}
Primijetimo da talas koji se prostire kao Gausov snop fazno kasni za odgovarajućim sfernim talasom, i za dovoljno veliko rastojanje pokupi fazno kašnjenje od:
\begin{equation*}
    \lim_{z \to \infty} \arctg\frac{z}{z_0} = \frac{\pi}{2}
\end{equation*}

\subsection{Poluprečnik krivine frontova}
Talasni frontovi su definisani jednačinom $\phi = C = const$:
\begin{align*}
     & kz + \frac{k\rho^2}{2R(z)} - \xi(z) = C             \\
     & k\left(z + \frac{\rho^2}{2R(z)}\right) = C + \xi(z) \\
     & z + \frac{\rho^2}{2R(z)} = \frac{C + \xi(z)}{k}
\end{align*}
Pošto se u dijelu prostora od interesa (u okolini ose, malo $\rho$) $R(z)$ i $\xi(z)$ malo mijenjaju u odnosu na $z$ i $\rho$, možemo aproksimirati $R(z) \approx const$ i $\xi(z) \approx const$ i napisati:
\begin{equation*}
    z + \frac{\rho^2}{2R} = const
\end{equation*}
što predstavlja jednačinu paraboloidne površi sa radijusom krivine R (analogno $y + \frac{x^2}{2R} = const$).

Primijetimo još da je $R(z \to 0) \to \infty$, što znači da se u okolini snopa Gausov snop ponaša kao ravanski talas, dok je $R(z \gg z_0) \approx z$, što znači da se za veliko $z$ ponaša kao sferni talas.

(Grafik R(z))

\subsection{Osobine Gausovog snopa na specijalnim mjestima}

\begin{enumerate}
    \item na mjestu spota:

          Snop se ponaša kao ravanski talas ($R \to \infty$).
          U blizini centra snopa, za $z \ll z_0, \rho \ll W_0$ je $W(z) \approx W_0$, odnosno iradijansa snopa $I \propto e^{-\frac{\rho^2}{W_0^2}} \approx 1$ je približno konstantna.

    \item na Rejlijevoj dužini ($z = z_0$):

          Intenzitet na osi jednak je polovini maksimuma.
          širina snopa je $\sqrt{2}$ puta veća od struka, a površina je veća 2 puta.
          Faza na osi je zakašnjena za $\frac{\pi}{4}$ u odnosu na ravanski talas.
          Talasni front je najzakrivljeniji sa radijusom $R_{min} = R(z_0) = 2z_0$.

    \item daleko od struka ($\rho < W_0, z \gg z_0$):

          Talas se približno ponaša kao sferni $R(z) \approx z$.
          Snop se širi linearno sa rastojanjem $W(z) \approx W_0 \frac{z}{z_0}$.
\end{enumerate}

\section{Karakterizacija Gausovog snopa}

Pored talasne dužine, za karakterizaciju različitih talasa potrebni su nam različiti parametri:
\begin{enumerate}
    \item za ravanski talas, potrebna su nam 2 parametra: pravac i kompleksna amplituda
    \item za sferni talas, potrebna su takođe 2 parametra: mjesto nastanka i kompleksna amplituda
    \item za Gausov snop, potrebna su nam 4 parametra: maksimalna amplituda $I_0$, položaj ose snopa, položaj struka i poluprečnik struka $W_0$ (ili Rejlijeva dužina $z_0$)
\end{enumerate}
\begin{tbox}
    Rejlijeva dužina ili zona označava prelazak iz polja bliske zone u polje daleke zone.
\end{tbox}
\begin{tbox}
    Za snop poznate amplitude i ose $q$ parametar određuje snop.
\end{tbox}

\subsection{Kvalitet Gausovog snopa}
Kvalitet realnog Gausovog snopa se određuje na osnovu odstupanja od idealnog, za koji (podsjetimo se) važi da je $2W_0 \cdot 2\theta_0 = \frac{4\lambda}{\pi}$.
Zbog toga se kao dobar pokazatelj kvaliteta snopa uzima parametar:
\begin{equation*}
    M^2 = \frac{2W_m2\theta_m}{\frac{4\lambda}{\pi}}
\end{equation*}
gdje su $W_m$ i $\theta_m$ izmjerene vrijednosti širine struka i divergencije realnog snopa, respektivno.
Ovaj izraz se koristi i za karakterizaciju svih snopova (ne samo Gausovog), zato što od svih njih Gausov ima najmanju divergenciju.
Ako uzmemo da je $W_m = W_0$, ostaje da je:
\begin{equation*}
    M^2 = \frac{\theta_m}{\theta_0} \geq 1 \text{ (divergencija snopa mora biti veća nego u idealnom slučaju)}
\end{equation*}
Na primjer, za HeNe lasere može važiti i $M < 1.1$, za jonske lasere je $M$ između 1.1 i 1.3, a za diodne između 1.1 i 1.7.
Za naš HeNe laser koji ćemo koristiti na laboratorijskim vježbama:
\begin{equation*}
    M^2 = \frac{\pi W_m \theta_m}{\lambda} = \pi \frac{\frac{0.63\si{\milli\meter}}{2}\cdot 1.3\si{\milli\radian}}{633\si{\nano\meter}} = 2
\end{equation*}

\subsection{Eksperimentalno određivanje parametara Gausovog snopa}

Da bi poznavali Gausov snop, potrebno je poznavati 5 parametara:
\begin{itemize}
    \item pravac prostiranja (osa snopa)
    \item snaga $P$
    \item širina struka $W_0$
    \item talasna dužina $\lambda$ ili Rejlijeva dužina $z_0$ ($z_0 = \frac{\pi W_0^2}{\lambda}$)
    \item položaj $z = 0$
\end{itemize}
Ravan $z = 0$ nije poznata unaprijed i zavisi od rezonatora - talasni frontovi se zbog graničnih uslova moraju prilagoditi ogledalima rezonatora.

Pozabavimo se određivanjem $W_0$.
To je moguće uraditi mjerenjem poluprečnika snopa $W_1 = W(z_1)$ i $W_2 = W(z_2)$ na dva proizvoljna položaja $z_1$ i $z_2$.
Razlika položaja $\Delta z$ ne zavisi od položaja koordinatnog početka (koji je nepoznat), pa uz pomoć vrijednosti $W_1$, $W_2$ i $\Delta z$ možemo odrediti $W_0$:
\begin{align*}
     & W = W_0\sqrt{1 + \left(\frac{z}{z_0}\right)^2}                                                                                                                   \\
     & \left(\frac{W}{W_0}\right)^2 = 1 + \left(\frac{z}{z_0}\right)^2                                                                                                  \\
     & \frac{W^2 - W_0^2}{W_0^2} z_0^2 = z^2                                                                                                                            \\
     & z = \frac{z_0}{W_0} \sqrt{W^2 - W_0^2} = \frac{\frac{\pi W_0^{\cancel2}}{\lambda}}{\cancel{W_0}} \sqrt{W^2 - W_0^2} = \frac{\pi W_0}{\lambda} \sqrt{W^2 - W_0^2} \\
     & z_2 - z_1 = \Delta z = \frac{\pi W_0}{\lambda} \left(\sqrt{W_2^2 - W_0^2} - \sqrt{W_1^2 - W_0^2}\right)                                                          \\
\end{align*}
Rješavanjem kvadratne jednačine po $W_0$ i biranjem pozitivnog rješenja dobija se:
\begin{equation*}
    W_0^2 = \frac{\frac{\lambda^2}{\pi^2}\Delta z^2\left(W_1^2 + W_2^2 + 2\sqrt{W_1^2W_2^2 - \frac{\lambda^2}{\pi^2}\Delta z^2}\right)}{\left(W_1^2 - W_2^2\right)^2 + 4\frac{\lambda^2}{\pi^2}\Delta z^2}
\end{equation*}
Kada je poznato $W_0$, možemo tu vrijednost ubaciti u $z_1 = \frac{z_0}{W_0}\sqrt{W_1^2 - W_0^2}$ i odrediti koliko je $z_1$, čime nam je poznat i položaj koordinatnog početka $z = 0$.

\section{Transmisija Gausovog snopa kroz optičke komponente}

Pri prolasku kroz optičke komponente rotaciono simetrične u odnosu na osu po kojoj se prostire snop, Gausov snop ostaje isti po formi, ukoliko se održi paraksijalna aproksimacija.
Ono što mu se mijenja jesu struk i širenje snopa.

\subsection{Prolazak Gausovog snopa kroz tanko sočivo}
Kompleksna transmitansa tankog sočiva žižne daljine $f$ je $t(\rho) = e^{j\frac{k\rho^2}{2f}}$ ($|t| = 1$, odnosno apsorpscija i refleksija su zanemarljive, dolazi samo do promjene faze).
Ako je incidentno električno polje $E$, polje pri izlasku iz sočiva iznosi:
\begin{equation*}
    E' = E \cdot t = \underbrace{\Psi_0' \frac{W_0}{W(z)} e^{-\frac{\rho^2}{W(z)^2}} e^{-jk\frac{\rho^2}{2R(z)} - j\xi - jkz}}_E \underbrace{e^{j\frac{k\rho^2}{2f}}}_t
\end{equation*}
Dakle, faktor $t$ mijenja talasni front na mjestu sočiva, dok je širina snopa očigledno ista $W' = W$.
Talasne dužina, a samim tim i talasni broj, se takođe ne mijenja.
Fazu transmitovanog talasa možemo zapisati kao fazu novog Gausovog snopa, u novom koordinatnom sistemu i novim parametrima $R'$, $W_0'$, $z_0'$ i $\xi'$:
\begin{equation*}
    k \frac{\rho^2}{2R} + kz + \xi - \frac{k\rho^2}{2f} = k \frac{\rho^2}{2R'} + kz' + \xi'
\end{equation*}
Imamo novu vrijednost $z_0'$, ali možemo da izaberemo novi koordinatni početak tako da je $kz + \xi = kz' + \xi'$ ($\xi = \arctg\frac{z}{z_0}$):
\begin{align*}
     & \cancel{k}\frac{\cancel{\rho^2}}{\cancel{2}R} - \cancel{k}\frac{\cancel{\rho^2}}{\cancel{2}f} = \cancel{k}\frac{\cancel{\rho^2}}{\cancel{2}R'} \\
     & \frac{1}{R} - \frac{1}{f} = \frac{1}{R'}
\end{align*}
Dakle, za snop koji je prošao kroz sočivo poznato nam je $W' = W$ i $R' = \frac{fR}{f - R}$.
\begin{tbox}
    Za divergirajući snop je $R > 0$, a za konvergirajući $R < 0$.
\end{tbox}
\begin{tbox}
    Za proizvoljan Gausov snop, čija su širina i poluprečnik poznati na nekoj koordinati $z$, moguće je odrediti poluprečnik i rastojanje struka od te tačke:
    \begin{align*}
         & W(z) = W_0 \sqrt{1 + \left(\frac{z}{z_0}\right)^2}                                                                                                                                                                                                                                                                                             \\
         & R(z) = z\left(1 + \left(\frac{z_0}{z}\right)^2\right)                                                                                                                                                                                                                                                                                          \\
         & \frac{W^2}{R} = \frac{W_0^2 \left(1 + \left(\frac{z}{z_0}\right)^2\right)}{z \left(1 + \left(\frac{z_0}{z}\right)^2\right)} = \frac{W_0^2 \frac{\cancel{z_0^2 + z^2}}{z_0^2}}{\cancel{z}\frac{\cancel{z^2 + z_0^2}}{z^{\cancel{2}}}} = W_0^2 \frac{z}{z_0^2} = W_0^2 \frac{z}{\frac{W_0^4 \pi^2}{\lambda^2}} = z \frac{\lambda^2}{\pi^2 W_0^2} \\
         & \implies z = \frac{\pi^2 W_0^2 W^2}{R \lambda^2}                                                                                                                                                                                                                                                                                               \\
         & R = \frac{z^2 + z_0^2}{z} = \frac{\frac{\pi^{\cancel{4} 2} W_0^4 W^4}{R^2 \lambda^{\cancel{4}2}} + \frac{W_0^4 \cancel{\pi^2}}{\cancel{\lambda^2}}}{\frac{\cancel{\pi^2} W_0^2 W^2}{R \cancel{\lambda^2}}}
    \end{align*}
    Odavde možemo izraziti $W_0$:
    \begin{align*}
         & \cancel{R} \cdot \frac{W_0^2 W^2}{\cancel{R}} = \frac{\pi^2 W_0^4 W^4}{R^2 \lambda ^2} + W_0^4 \bigg/ :W_0^2 \\
         & W^2 = \frac{\pi^2 W_0^2 W^4}{R^2 \lambda^2} + W_0^2 = W_0^2 \left(1 + \frac{\pi^2 W^4}{R^2 \lambda^2}\right) \\
         & W_0^2 = \frac{W^2}{1 + \frac{\pi^2 W^4}{R^2 \lambda^2}}
    \end{align*}
    \begin{empheq}[box=\fbox]{equation*}
        W_0 = \frac{W}{\sqrt{1 + \left(\frac{\pi W^2}{R \lambda}\right)^2}}
    \end{empheq}
    \begin{align*}
        z = \frac{\pi^2 W^2}{R \lambda^2} W_0^2 = \frac{\pi^2 W^2}{R \lambda^2} \frac{W^2}{1 + \left(\frac{\pi W^2}{R \lambda}\right)^2}
    \end{align*}
\end{tbox}

\end{document}